\documentclass[11pt,a4paper]{article}

\usepackage[T1]{fontenc}
\usepackage[utf8]{inputenc}
\usepackage[italian]{babel}
\usepackage{lmodern}
\usepackage{microtype}
\usepackage[a4paper,margin=2.25cm,headheight=15pt]{geometry}
\usepackage{booktabs,longtable,tabularx,array}
\usepackage{ragged2e}
\usepackage{enumitem}
\usepackage{xcolor}
\usepackage[most]{tcolorbox}
\usepackage{fancyhdr}
\usepackage{hyperref}
\usepackage{xurl}
\usepackage{listings}
\usepackage{tikz}
\usetikzlibrary{arrows.meta,positioning,calc,fit}
\usepackage{titlesec}
\usepackage{amsmath,amssymb}

\definecolor{ddtadark}{RGB}{28,38,52}
\definecolor{ddtagray}{RGB}{244,246,248}
\definecolor{ddtablue}{RGB}{42,88,135}
\definecolor{ddtagreen}{RGB}{48,111,78}
\definecolor{ddtaorange}{RGB}{148,91,25}
\definecolor{ddtared}{RGB}{140,48,48}
\definecolor{ddtapurple}{RGB}{92,63,122}

\hypersetup{
  colorlinks=true,
  linkcolor=ddtablue,
  urlcolor=ddtablue,
  citecolor=ddtablue,
  pdftitle={DDTA Documentation and Base Analysis Authoring Guide R2},
  pdfauthor={Documentation-Driven Threat Analysis research workspace}
}

\pagestyle{fancy}
\fancyhf{}
\lhead{DDTA -- Authoring Guide R2}
\rhead{Documentation method}
\cfoot{\thepage}

\titleformat{\section}{\Large\bfseries\color{ddtadark}}{\thesection}{0.7em}{}
\titleformat{\subsection}{\large\bfseries\color{ddtadark}}{\thesubsection}{0.7em}{}
\titleformat{\subsubsection}{\normalsize\bfseries\color{ddtadark}}{\thesubsubsection}{0.7em}{}

\setlist[itemize]{leftmargin=1.55em,itemsep=0.24em,topsep=0.35em}
\setlist[enumerate]{leftmargin=1.75em,itemsep=0.24em,topsep=0.35em}
\setlength{\parskip}{0.35em}
\setlength{\parindent}{0pt}
\setlength{\emergencystretch}{3em}

\lstset{
  basicstyle=\ttfamily\small,
  columns=fullflexible,
  keepspaces=true,
  frame=single,
  framerule=0.3pt,
  rulecolor=\color{ddtablue},
  backgroundcolor=\color{ddtagray},
  xleftmargin=0.6em,
  xrightmargin=0.6em,
  aboveskip=0.5em,
  belowskip=0.5em,
  breaklines=true,
  showstringspaces=false
}

\newtcolorbox{statusclosed}[1][CLOSED]{enhanced,breakable,colback=ddtagray,colframe=ddtagreen,title={#1},fonttitle=\bfseries}
\newtcolorbox{statusworking}[1][R24 WORKING]{enhanced,breakable,colback=white,colframe=ddtaorange,title={#1},fonttitle=\bfseries}
\newtcolorbox{statusopen}[1][OPEN / DEFERRED]{enhanced,breakable,colback=white,colframe=ddtapurple,title={#1},fonttitle=\bfseries}
\newtcolorbox{goodbox}[1][CORRETTO / PREFERIBILE]{enhanced,breakable,colback=white,colframe=ddtagreen,title={#1},fonttitle=\bfseries}
\newtcolorbox{badbox}[1][DA EVITARE]{enhanced,breakable,colback=white,colframe=ddtared,title={#1},fonttitle=\bfseries}
\newtcolorbox{reviewbox}[1][REVIEW]{enhanced,breakable,colback=white,colframe=ddtaorange,title={#1},fonttitle=\bfseries}
\newtcolorbox{examplebox}[1][ESEMPIO FACIAL ACCESS]{enhanced,breakable,colback=ddtagray,colframe=ddtablue,title={#1},fonttitle=\bfseries}

\newcolumntype{Y}{>{\RaggedRight\arraybackslash}X}
\newcolumntype{L}[1]{>{\RaggedRight\arraybackslash}p{#1}}

\newcommand{\MR}{\textit{MacroRequirement}}
\newcommand{\Decision}{\textit{Decision}}
\newcommand{\FR}{\textit{FunctionalRequirement}}
\newcommand{\Req}{\textit{Requirement}}
\newcommand{\SR}{\textit{SpecializedRequirement}}
\newcommand{\SecR}{\textit{SecurityRequirement}}
\newcommand{\BA}{\textit{Base Analysis}}
\newcommand{\code}[1]{\texttt{#1}}

\begin{document}

\begin{titlepage}
\centering
\vspace*{1.5cm}
{\Large Documentation-Driven Threat Analysis\par}
\vspace{0.8cm}
{\Huge\bfseries Documentation and Base Analysis Authoring Guide\par}
\vspace{0.25cm}
{\LARGE Revision 2 -- documentation method\par}
\vspace{0.9cm}
{\large Guida operativa dalla problem framing al handoff verso Base Analysis\par}
\vfill
\begin{minipage}{0.86\textwidth}
\small
\textbf{Stato.} R24 DRAFT FOR REVIEW. La guida e' derivata dal checkpoint documentale R24 e conserva esplicitamente la distinzione fra regole CLOSED, refinement R24 WORKING e questioni OPEN/DEFERRED.

\medskip
\textbf{Repository baseline di riferimento.} \code{7aeb3a1e0373357fa31c35110d9160de3e2e7d79}.

\medskip
\textbf{Esempio ricorrente.} Il caso Facial Access viene usato come esempio didattico. Gli artifact R1 sono SUPERSEDED e servono solo come regression evidence; il candidate R2 e' EXPERIMENTAL\_NON\_CANONICAL e non costituisce ancora una baseline governata per una BA accettata.
\end{minipage}
\vfill
{\large 23 agosto 2026\par}
\end{titlepage}

\pagenumbering{roman}
\tableofcontents
\newpage
\pagenumbering{arabic}

\section{Come usare questa guida}

Questa guida non sostituisce la documentazione di progetto con un metamodel pesante. Il principio operativo e':

\begin{statusclosed}[Principio -- minimum sufficient governed meaning]
Documentare al livello corrente la minima conoscenza governata sufficiente a rendere stabile il significato, verificabile la responsabilita' e possibile una successiva normalizzazione in Base Analysis senza inventare dettagli.
\end{statusclosed}

La sequenza di lavoro e':

\begin{center}
\begin{tikzpicture}[
  node distance=6mm,
  box/.style={draw=ddtadark,rounded corners,align=center,inner sep=5pt,minimum width=62mm},
  arr/.style={-{Latex[length=2.3mm]},thick}
]
\node[box] (a) {0. Authority gate};
\node[box,below=of a] (p) {1. Project problem framing};
\node[box,below=of p] (m) {2. MacroRequirement};
\node[box,below=of m] (s) {3. Semantic-sufficiency gate};
\node[box,below=of s] (d) {4. Decision};
\node[box,below=of d] (f) {5. FunctionalRequirement};
\node[box,below=of f] (r) {6. Requirement coherent-unit / split};
\node[box,below=of r] (sr) {7. SpecializedRequirement};
\node[box,below=of sr] (sec) {8. SecurityRequirement};
\node[box,below=of sec] (svc) {9. Cross-MR / consumed-service boundary};
\node[box,below=of svc] (t) {10. Canonical terminology};
\node[box,below=of t] (reg) {11. Downstream semantic propagation};
\node[box,below=of reg] (g) {12. Completeness / promotion gate};
\node[box,below=of g] (ba) {13. Handoff to Base Analysis};
\draw[arr] (a)--(p);\draw[arr] (p)--(m);\draw[arr] (m)--(s);\draw[arr] (s)--(d);\draw[arr] (d)--(f);\draw[arr] (f)--(r);\draw[arr] (r)--(sr);\draw[arr] (sr)--(sec);\draw[arr] (sec)--(svc);\draw[arr] (svc)--(t);\draw[arr] (t)--(reg);\draw[arr] (reg)--(g);\draw[arr] (g)--(ba);
\end{tikzpicture}
\end{center}

Non ogni step crea un nuovo tipo di documento. Alcuni sono gate o review step. In particolare semantic sufficiency, Macro Project Map review, downstream regression e service-guarantee review sono discipline di controllo, non nuove classi L1.

\subsection{Legenda epistemica}

\begin{longtable}{L{0.25\textwidth}L{0.64\textwidth}}
\toprule
\textbf{Etichetta} & \textbf{Significato} \\
\midrule\endhead
CLOSED & Regola gia' chiusa nel proprio scope e portata avanti nel baseline corrente. \\
R24 WORKING & Refinement supportato da pressure test ma non ancora generalizzato come closure universale. \\
OPEN / DEFERRED & Questione preservata esplicitamente; non deve essere chiusa per comodita' narrativa o di tooling. \\
\bottomrule
\end{longtable}

\section{Regole trasversali prima dei singoli documenti}

\subsection{Recency non e' authority}

\begin{statusclosed}
Prima di usare un artifact chiedere sempre quale baseline e quale source set sono autorizzati per l'operazione corrente. Un file piu' nuovo non diventa governato per cronologia.
\end{statusclosed}

\begin{badbox}
\begin{lstlisting}
"Questo file e' il piu' recente, quindi e' il requisito corrente."
\end{lstlisting}
\end{badbox}

\begin{goodbox}
\begin{lstlisting}
"Questo artifact e' candidate / working evidence.
La baseline governata si determina tramite il registry di authority."
\end{lstlisting}
\end{goodbox}

\subsection{Documentation authority e analysis feedback}

La direzione di authority e':

\begin{lstlisting}
governed documentation
    -> Base Analysis
    -> downstream analysis / diagnostic
    -> correction candidate
    -> governed review
    -> updated documentation
    -> rebuilt Base Analysis
\end{lstlisting}

Una BA o un'analisi possono esporre un problema, ma non possono riparare silenziosamente il significato del progetto.

\section{Step 1 -- Project problem framing}

\begin{statusclosed}
\textbf{Domanda fondamentale:} se rimuoviamo le scelte di soluzione correnti, quale problema generale o classe di problemi deve affrontare il progetto, e che cosa e' dentro o fuori dal boundary?
\end{statusclosed}

La problem framing e' una precondizione di metodo, non un ulteriore document type L1.

\subsection{Checklist}

\begin{itemize}
\item eliminare provider, device, database, framework, algoritmo e protocollo quando non sono intrinseci al problema;
\item dichiarare le responsabilita' importanti fuori scope;
\item usare il framing per verificare copertura e overlap degli MR;
\item non duplicare meccanicamente lo stesso framing in ogni MR.
\end{itemize}

\begin{badbox}[DA EVITARE -- framing gia' soluzione]
\begin{lstlisting}
Il progetto deve usare una telecamera IP Ethernet e un PC locale
per riconoscere il volto e aprire il tornello.
\end{lstlisting}
Problema: camera, Ethernet, PC e apertura del tornello sono gia' scelte o realizzazioni. Il framing non lascia visibile il problema indipendentemente dalla soluzione.
\end{badbox}

\begin{goodbox}[PREFERIBILE -- framing del problema]
\begin{lstlisting}
Il progetto deve controllare l'accesso a un'area riservata,
distinguendo la persona presente, il relativo stato di autorizzazione
e la decisione sul singolo tentativo di accesso.
\end{lstlisting}
La formulazione sopravvive alla sostituzione della camera, del riconoscimento facciale, della rete o del meccanismo fisico di apertura.
\end{goodbox}

\section{Step 2 -- MacroRequirement}

\begin{statusclosed}
\textbf{Domanda fondamentale:} nel problema gia' inquadrato, quale responsabilita' macro stabile stiamo governando, quale risultato o valore macro contribuisce, perche' conta e chi coinvolge?
\end{statusclosed}

\subsection{Corpo minimo}

\begin{lstlisting}
Title
Intent
Context
Stakeholders
Scope
Assumptions/Constraints [solo se branch-wide]
dependsOn MR [quando genuino]
\end{lstlisting}

\subsection{Review questions}

\begin{itemize}
\item esiste una sola macro-responsabilita' coerente?
\item resta significativa cambiando tecnologia/architettura?
\item contiene gia' un comportamento atomico da FR?
\item contiene gia' una scelta di soluzione da Decision?
\item una dependency e' confusa con containment?
\item Intent + Context + Scope determinano un solo significato stabile?
\end{itemize}

\subsection{Esempio reale: MR-0003 prima e dopo la semantic review}

Il superseded R1 recitava, in sintesi:

\begin{examplebox}[REGRESSION EVIDENCE -- formulazione R1]
\begin{lstlisting}
MR-0003 -- Verifica dell'identita' al punto di accesso

Intent:
stabilire se la persona presente corrisponde a un'identita' governata.
\end{lstlisting}
\end{examplebox}

Questa formulazione ammetteva due letture materialmente diverse:

\begin{lstlisting}
A. l'identita' specifica non e' ancora nota;
   la responsabilita' determina QUALE GovernedIdentity corrisponde.

B. una GovernedIdentity specifica e' gia' selezionata;
   la responsabilita' determina SE la persona corrisponde a quella identita'.
\end{lstlisting}

Non era quindi ``sbagliata'' in senso sintattico: era semanticamente insufficiente per distinguere due responsabilita' diverse.

\begin{goodbox}[CANDIDATE R2 -- responsabilita' resa stabile]
\begin{lstlisting}
MR-0003 -- Determinazione dell'identita' al punto di accesso

Quando questa responsabilita' inizia, la persona e' presente
ma la specifica identita' governata che le corrisponde
non e' ancora determinata.

La responsabilita' determina quale identita' governata
corrisponde alla persona presente.
\end{lstlisting}
\end{goodbox}

\begin{badbox}[DA EVITARE -- MR che assorbe la realizzazione]
\begin{lstlisting}
MR-0003 -- Riconoscimento facciale tramite CameraSubsystem
su Ethernet verso RecognitionProcessor.
\end{lstlisting}
Qui strategia, placement e comunicazione sono anticipati nel MacroRequirement. La sostituzione del mezzo o del componente costringerebbe a riscrivere una responsabilita' che dovrebbe restare stabile.
\end{badbox}

\subsection{Macro Project Map / isolated-MR review}

\begin{statusworking}[R24 / S1 CANDIDATE HEURISTIC]
Un MR isolato non e' invalido. Un nodo isolato attiva review, soprattutto quando il contenuto sembra una raccolta di proprieta' trasversali su comportamenti posseduti da altri branch.
\end{statusworking}

Domande:

\begin{enumerate}
\item Rimuovendo l'MR scompare una vera macro-capability/responsabilita'?
\item Oppure restano le capability e scompaiono principalmente proprieta' specializzate?
\item I discendenti plausibili sarebbero soprattutto security/privacy/performance/compliance requirements su FR altrui?
\item Manca una vera relazione \code{dependsOn}?
\item L'MR e' stato creato solo per dare una ``casa'' a una metodologia o concern category?
\end{enumerate}

\section{Step 3 -- Semantic-sufficiency gate}

\begin{statusworking}
Questo step e' un refinement R24 sotto pressure test. Le dimensioni A--F sono dimensioni di osservazione, non campi obbligatori del metamodel.
\end{statusworking}

\textbf{Trigger question:}

\begin{quote}
Ignorando titolo e nomi suggestivi, il testo governato determina una sola responsabilita'/commitment/obligation stabile, oppure possono sopravvivere letture materialmente diverse?
\end{quote}

\subsection{Dimensioni di osservazione}

\begin{lstlisting}
A  context / boundary
B  participants / domains / relations
C  initial / known / observable state
D  events / activities / state transitions
E  required effect / outcome / responsibility
F  domain properties / constraints
\end{lstlisting}

\subsection{Procedura minima}

\begin{enumerate}
\item ricostruire il significato usando neutral labels quando utile;
\item formare le piu' piccole letture materialmente diverse;
\item individuare la proposizione minima che le separa;
\item chiedere se la differenza cambia responsabilita', relazione, stato, outcome o downstream structure;
\item cercare evidenza governata;
\item classificare la proposizione come \code{AFFIRMED}, \code{DENIED}, \code{NOT SPECIFIED} o \code{CONFLICTING};
\item correggere il livello che possiede semanticamente il significato.
\end{enumerate}

\begin{examplebox}[MR-0003 -- critical proposition]
\begin{lstlisting}
Una specifica GovernedIdentity e' gia' disponibile/selezionata
come riferimento prima che inizi la responsabilita'.
\end{lstlisting}
Nel caso Facial Access la clarification di progetto ha reso la proposizione \code{DENIED}; la responsabilita' deve determinare quale identita' corrisponde.
\end{examplebox}

\begin{badbox}[DA EVITARE -- inference che chiude il gap]
\begin{lstlisting}
"Si chiama verifica, quindi certamente esiste gia'
un'identita' specifica da verificare."
\end{lstlisting}
Titolo, pattern, esperienza di dominio o LLM possono suggerire una lettura, non diventano project authority.
\end{badbox}

\begin{goodbox}[CORRETTO -- mantenere l'incertezza]
\begin{lstlisting}
Evidence: NOT SPECIFIED
Action: formulate discriminator question / return to semantic owner.
\end{lstlisting}
\end{goodbox}

\section{Step 4 -- Decision}

\begin{statusclosed}[CLOSED CORE / FORWARD WORDING REFINED]
\textbf{Domanda fondamentale:} quale \emph{material project commitment} restringe questo MR, perche' il progetto assume questa posizione e quali conseguenze materiali seguono?
\end{statusclosed}

Usare \emph{material project commitment} evita l'equivoco che ogni Decision debba rappresentare una libera scelta fra molte alternative. Il commitment puo' essere \code{SELECTABLE}, \code{NECESSITY-CONSTRAINED} o \code{DEFAULT}, ma non puo' essere un wrapper vuoto.

\subsection{Corpo minimo}

\begin{lstlisting}
Title
Context
Decision
Consequences
\end{lstlisting}

\subsection{Ownership test}

Per ogni Decision chiedere anche:

\begin{quote}
Questa proposizione e' una scelta/necessita' adottata per realizzare questo MR, oppure descrive semplicemente una proprieta' di qualcosa che il branch consuma?
\end{quote}

\begin{goodbox}[CORRETTO -- Decision direttamente pertinente a MR-0003]
\begin{lstlisting}
D-3.1 -- Strategia di determinazione dell'identita'
mediante riconoscimento facciale

Decision:
Il progetto usa riconoscimento facciale come strategia
per determinare quale identita' governata corrisponde
alla persona presente.
\end{lstlisting}
La Decision restringe direttamente la realizzazione di MR-0003 senza cambiare la responsabilita' macro.
\end{goodbox}

\begin{badbox}[DA EVITARE -- proprieta' infrastrutturale presentata automaticamente come Decision del MR]
\begin{lstlisting}
D-3.x -- La rete e' Ethernet cablata.
\end{lstlisting}
Se ``Ethernet cablata'' e' solo una proprieta' del servizio di comunicazione disponibile e potrebbe servire anche altri branch, non e' automaticamente una Decision posseduta da MR-0003. Serve prima il test di ownership/placement.
\end{badbox}

\begin{badbox}[DA EVITARE -- implementation detail anticipato]
\begin{lstlisting}
D-3.1 -- Usare il modello FaceNet vX, soglia 0.83,
telecamera modello Y e TLS 1.3.
\end{lstlisting}
La Decision miscela strategia, algoritmo, threshold, dispositivo e protezione del trasporto. Questi commitment possono avere ownership e lifecycle differenti.
\end{badbox}

\section{Step 5 -- FunctionalRequirement}

\begin{statusclosed}
\textbf{Domanda fondamentale:} data la Decision parent e gli input/condizioni applicabili, che cosa MUST fare il subject/capability, su quale significato governato, e quale risultato osservabile o failure behavior MUST seguire?
\end{statusclosed}

\subsection{Review checklist}

\begin{itemize}
\item comportamento operativo;
\item independently assessable;
\item esattamente una parent Decision;
\item parent Decision pertinente;
\item non ripetere la Decision;
\item una sola unita' funzionale coerente;
\item separare comportamenti indipendentemente modificabili;
\item implementation independence;
\item riuso delle identita'/riferimenti governati;
\item testability senza inventare significato.
\end{itemize}

\begin{goodbox}[CANDIDATE R2 -- FR-3.2.1]
\begin{lstlisting}
Quando viene eseguita la determinazione dell'identita',
la capability di riconoscimento MUST produrre un esito governato
e MUST mantenere distinguibili almeno:
- determinazione riuscita;
- determinazione negativa;
- esito non conclusivo.

Quando la determinazione riesce, MUST rendere disponibile
la specifica identita' governata determinata come corrispondente.
\end{lstlisting}
La clausola governa il comportamento e il risultato, senza decidere authorization/access e senza imporre il meccanismo interno.
\end{goodbox}

\begin{badbox}[DA EVITARE -- realism non governato]
\begin{lstlisting}
Il RecognitionProcessor MUST identificare la persona
quando confidence > 0.82 e MUST usare top-1 matching 1:N.
\end{lstlisting}
Nel baseline corrente confidence, score, threshold, ranking e semantica 1:N non sono governati. Introdurli per ``realismo'' rende la documentazione meno stabile e inventa project meaning.
\end{badbox}

\begin{badbox}[DA EVITARE -- FR che invade authorization/access decision]
\begin{lstlisting}
Se il volto viene riconosciuto, il sistema MUST autorizzare
l'accesso e aprire il varco.
\end{lstlisting}
Il riconoscimento determina l'identita'; autorizzazione e decisione del singolo tentativo appartengono ad altre responsabilita'.
\end{badbox}

\subsection{Allowed domain vs decision rule}

\begin{goodbox}
\begin{lstlisting}
Outcome MUST be one of:
  SUCCESS
  NEGATIVE
  INCONCLUSIVE
\end{lstlisting}
Questo governa un dominio di risultati.
\end{goodbox}

\begin{badbox}
\begin{lstlisting}
IF score >= 0.82 THEN SUCCESS
ELSE IF score < 0.40 THEN NEGATIVE
ELSE INCONCLUSIVE
\end{lstlisting}
Se soglie e mapping non sono governati, questa apparente testabilita' e' semantic invention.
\end{badbox}

\section{Step 6 -- Common Requirement coherent-unit / split}

\begin{statusclosed}[S1.5-A CLOSED]
Un Requirement e' esso stesso una obbligazione normativa governata e possiede una o piu' normative clauses.
\end{statusclosed}

\begin{lstlisting}
Requirement [abstract]
    extends GovernedDocument
    normativeClause : NormativeClause [1..*]
\end{lstlisting}

Regola:

\begin{lstlisting}
1 Requirement != 1 sentence
1 Requirement  = 1 coherent normative obligation
\end{lstlisting}

\textbf{Split test:} una parte puo' essere introdotta, modificata, ritirata o valutata indipendentemente senza cambiare l'identita' normativa delle altre? Se si', valutare Requirement separati.

\begin{badbox}[DA EVITARE -- generic secure requirement]
\begin{lstlisting}
RecognitionCapture MUST be secure, confidential,
integrity-protected and from an authorized source.
\end{lstlisting}
Confidentiality, integrity e authorized provenance possono evolvere ed essere valutate indipendentemente. Un unico requisito rende difficile ownership, change e verifica.
\end{badbox}

\begin{goodbox}[PREFERIBILE -- split per obbligazioni indipendenti]
\begin{lstlisting}
SEC-C  Confidentiality
SEC-I  Integrity
SEC-P  Authorized provenance
\end{lstlisting}
Ogni SecurityRequirement conserva una sola obbligazione coerente e una singola protected security property.
\end{goodbox}

\section{Step 7 -- SpecializedRequirement}

\begin{statusclosed}[S1 CLOSED + S1.5 COMMON ABSTRACTION]
\textbf{Domanda fondamentale:} quale obbligazione concern-specific aggiuntiva deve valere insieme al parent FR perche' la funzione sia considerata soddisfatta, senza ridefinire la funzione ne' scegliere come realizzarla?
\end{statusclosed}

Checks:

\begin{itemize}
\item esattamente un parent FR;
\item strengthening, non replacement;
\item removal test: tolto lo SR, la funzione ordinaria resta riconoscibile;
\item ordinary functional correctness non duplicata;
\item obbligazioni indipendenti separate;
\item subordinate positive behavior puo' esistere, ma resta subordinato al concern;
\item nessun control/mechanism obbligato senza Decision pertinente.
\end{itemize}

\begin{goodbox}
\begin{lstlisting}
Parent FR:
CameraSubsystem MUST deliver RecognitionCapture
al RecognitionProcessor corretto.

Specialized obligation:
Durante tale delivery, una capture modificata senza autorizzazione
MUST NOT poter essere accettata come valida.
\end{lstlisting}
Il secondo requisito restringe la soddisfazione del delivery ma non ridefinisce la funzione di delivery.
\end{goodbox}

\begin{badbox}
\begin{lstlisting}
SpecializedRequirement:
CameraSubsystem MUST deliver RecognitionCapture.
\end{lstlisting}
Questa clausola duplica la correttezza funzionale del parent FR invece di specializzarla.
\end{badbox}

\section{Step 8 -- SecurityRequirement}

\begin{statusclosed}[S2 CLOSED]
Un SecurityRequirement e' uno SpecializedRequirement che preserva una singola proprieta' di sicurezza governata rilevante per il parent FR e rende identificabile nelle normative clauses il failure mode che ne costituisce perdita/violazione.
\end{statusclosed}

\begin{lstlisting}
SecurityRequirement
    extends SpecializedRequirement
    protectedSecurityProperty : SecurityProperty [1]
\end{lstlisting}

\begin{goodbox}[CONFIDENTIALITY -- forma mechanism-neutral]
\begin{lstlisting}
Durante il delivery di RecognitionCapture,
il contenuto biometrico MUST NOT diventare intelligibile
a soggetti non autorizzati prima della sua accettazione
da parte del RecognitionProcessor.
\end{lstlisting}
Failure mode: il contenuto diventa intelligibile a un soggetto non autorizzato durante il comportamento protetto.
\end{goodbox}

\begin{badbox}[DA EVITARE -- control scambiato per requirement]
\begin{lstlisting}
RecognitionCapture MUST essere protetta con TLS 1.3 e mTLS.
\end{lstlisting}
TLS/mTLS sono possibili scelte di realization. Il SecurityRequirement deve prima dire quale proprieta' va preservata e quale failure mode e' inaccettabile.
\end{badbox}

\begin{badbox}[DA EVITARE -- causa scambiata per requisito]
\begin{lstlisting}
Il sistema MUST impedire gli attacchi man-in-the-middle.
\end{lstlisting}
Un attacco e' una possibile causa. Il requisito deve governare l'obbligazione sul comportamento/proprieta' del parent FR; la tecnica di analysis puo' poi collegare cause e failure mode.
\end{badbox}

\section{Step 9 -- Cross-MR e consumed-service boundary}

\subsection{Consumption vs ownership}

\begin{statusclosed}
Una branch puo' consumare una capability/servizio senza acquisirne ownership, senza un secondo parent e senza importare automaticamente la realizzazione interna del servizio.
\end{statusclosed}

\begin{lstlisting}
consumer branch
    -> consumes Service S

consumption != ownership
consumption != second parent
shared service != shared FR ownership
\end{lstlisting}

\begin{examplebox}[Camera -> RecognitionProcessor]
La separazione fra acquisizione in \code{CameraSubsystem} e riconoscimento in \code{RecognitionProcessor} crea la necessita' di rendere disponibile \code{RecognitionCapture} al processor. Questa interazione puo' consumare un servizio di comunicazione senza trasformare MR-0003 nel proprietario dell'infrastruttura di rete.
\end{examplebox}

\begin{badbox}[DA EVITARE -- importare la rete nel MR]
\begin{lstlisting}
MR-0003 deve fornire Ethernet, switch, VLAN e routing
per il trasferimento della capture.
\end{lstlisting}
Se la comunicazione e' una capability consumata, il consumer non possiede automaticamente questi dettagli.
\end{badbox}

\subsection{Required property vs governed service guarantee}

\begin{statusworking}[R24 WORKING INTEGRATION CANDIDATE]
Quando il consumer fa affidamento su una proprieta' P del servizio, confrontare la proprieta' richiesta con le sole garanzie governate del servizio. L'assenza di una garanzia non genera automaticamente un requisito o un meccanismo.
\end{statusworking}

\begin{lstlisting}
consumer requires property P
        |
        v
does governed service evidence guarantee P?
        |
        +-- AFFIRMED
        |      -> consumer may rely on the governed guarantee
        |
        `-- DENIED / NOT SPECIFIED / CONFLICTING
               -> coverage / assurance question remains
               -> determine semantic owner
               -> only then create/change an artifact if justified
\end{lstlisting}

\begin{goodbox}[CORRETTO -- informazione minima oggi]
\begin{lstlisting}
MR-0003 realization consumes CommunicationService
for RecognitionCapture delivery.

Service confidentiality guarantee: NOT SPECIFIED.
\end{lstlisting}
Questo consente di rimandare i dettagli di rete senza assumere proprieta' inesistenti.
\end{goodbox}

\begin{badbox}[DA EVITARE -- inferenza automatica]
\begin{lstlisting}
CommunicationService confidentiality = NOT SPECIFIED
therefore
MR-0003 MUST use TLS.
\end{lstlisting}
Il passaggio e' illegittimo: \code{NOT SPECIFIED} non significa ``assenza'', e anche un coverage gap richiede prima ownership/placement review.
\end{badbox}

\begin{reviewbox}[Regola pratica]
Non scrivere ``il canale e' sicuro''. Scrivere, quando serve, quali proprieta' il consumer richiede e quali proprieta' il servizio governa come garantite. La stessa communication capability puo' servire Camera->PC e PC->turnstile con requisiti diversi sui due flussi.
\end{reviewbox}

\section{Step 10 -- Canonical terminology e review bindings}

\begin{statusclosed}[BA5 DISCIPLINE / AUTHORING APPLICATION]
La prosa puo' restare naturale; i binding semanticamente operativi devono usare identita' e termini governati stabili. Grammatical variation e' ammessa, uncontrolled synonymy no.
\end{statusclosed}

\begin{goodbox}
\begin{lstlisting}
GovernedIdentity
GovernedIdentity determinata
alla GovernedIdentity corrispondente
\end{lstlisting}
Le flessioni grammaticali cambiano la superficie, non la semantic identity.
\end{goodbox}

\begin{badbox}
\begin{lstlisting}
GovernedIdentity
AuthorizedPerson
KnownUser
RecognizedProfile
\end{lstlisting}
Se questi termini indicano lo stesso significato senza una regola governata, introducono synonym drift e rendono piu' difficile capire se esistano concetti diversi.
\end{badbox}

\subsection{Temporary review markers}

Durante review possono essere usati marker come:

\begin{lstlisting}
*<RESPONSIBILITY_A>
#<GOVERNED_IDENTITY>
@<produce>
\end{lstlisting}

Devono restare visibili finche' servono alla review e devono essere rimossi nella prosa finale di progetto.

\section{Step 11 -- Downstream semantic propagation}

\begin{statusworking}[R24 WORKING -- ONE REAL DOWNSTREAM CASE]
Una correzione upstream richiede review dei discendenti non solo per contraddizione diretta, ma anche per perdita di una distinzione governata, precondizioni silenziose o ownership mal posizionata.
\end{statusworking}

\subsection{Facial Access propagation example}

Dopo la correzione di MR-0003, il downstream corretto deve preservare:

\begin{lstlisting}
entry:
  person present
  specific GovernedIdentity not yet determined

responsibility:
  determine which GovernedIdentity corresponds

success outcome:
  make that specific GovernedIdentity available
\end{lstlisting}

\begin{goodbox}[D-3.2 / FR-3.2.1 -- distinction preserved]
La Decision e l'FR candidate R2 dicono esplicitamente che il riconoscimento rende disponibile la specifica identita' determinata e non decide autorizzazione o accesso.
\end{goodbox}

\begin{badbox}[DA EVITARE -- child non contraddittorio ma lossy]
\begin{lstlisting}
Recognition MUST compare the person with GovernedIdentity
and output positive/negative.
\end{lstlisting}
La frase puo' non contraddire esplicitamente MR-0003, ma reintroduce la possibilita' che una specifica identita' sia gia' input alla responsabilita'.
\end{badbox}

Quando possibile, riusare la \emph{smallest critical proposition} come downstream regression question.

\section{Step 12 -- Documentation completeness e promotion gate}

\begin{statusclosed}[ACTIVE GOVERNANCE DISCIPLINE]
Una candidate baseline non diventa primary BA source perche' appare coerente o piu' recente. Serve una promozione esplicita dopo la composizione e review del successore completo.
\end{statusclosed}

Prima della promozione chiedere:

\begin{itemize}
\item tutti i branch richiesti sono assemblati coerentemente?
\item le correzioni upstream sono state propagate/reviewed?
\item i \code{NOT SPECIFIED} noti restano intenzionalmente aperti?
\item i documenti finali sono privi di marker temporanei e research commentary?
\item l'authority registry riflette il nuovo stato?
\end{itemize}

\begin{badbox}
\begin{lstlisting}
"Il candidate R2 e' scritto meglio, quindi lo usiamo per la BA corrente."
\end{lstlisting}
Una candidate non e' current authority senza promotion.
\end{badbox}

\section{Step 13 -- Handoff a Base Analysis}

\begin{statusclosed}
Documentation remains project authority. Base Analysis normalizza solo il significato gia' governato e non anticipa il dettaglio che il progetto ha intenzionalmente lasciato aperto.
\end{statusclosed}

Domanda di handoff:

\begin{quote}
Esiste sufficiente significato governato per derivare la minima Base Analysis senza scegliere interpretazioni non supportate o importare dettagli di livello inferiore?
\end{quote}

Se no, tornare al semantic owner nel document layer.

\begin{examplebox}[MR-0003 -- BA boundary]
Una futura BA derivata dal successor governato non deve richiedere una specifica \code{GovernedIdentity} come input preesistente della responsabilita' di determinazione. Se la BA la richiede, si attiva una semantic regression review: potrebbe essere un errore di derivazione, una perdita espressiva della BA o un problema documentale, ma non una nuova project truth.
\end{examplebox}

\section{Step 14 -- BA / analysis feedback senza authority inversion}

Un mismatch BA non dimostra automaticamente un difetto documentale. Le diagnosi minime sono:

\begin{lstlisting}
documentation ambiguity
BA derivation error
missing governed fact
BA expressiveness problem
\end{lstlisting}

Il loop corretto e':

\begin{lstlisting}
governed documentation
    -> minimum justified BA
    -> semantic / structural regression
    -> diagnostic
    -> correction candidate
    -> governed review
    -> updated documentation
    -> rebuilt BA
\end{lstlisting}

\section{Checklist compatta documento per documento}

\begin{longtable}{L{0.18\textwidth}L{0.73\textwidth}}
\toprule
\textbf{Step} & \textbf{Domande da non saltare} \\
\midrule\endhead
Authority & Quale baseline e' autorizzata? Sto usando recency come authority? \\
Framing & Il problema resta valido rimuovendo la soluzione corrente? Boundary chiaro? \\
MR & Una macro-responsabilita'? Significato stabile? Niente scelta tecnica o FR atomico? \\
Semantic sufficiency & Esistono letture materialmente diverse? Qual e' la smallest critical proposition? Evidenza AFFIRMED/DENIED/NOT SPECIFIED/CONFLICTING? \\
Decision & Quale material commitment restringe il MR? E' davvero posseduto da questo MR? Conseguenze materiali? \\
FR & Quale comportamento MUST? Risultato osservabile? Una coherent functional unit? Testabile senza inventare significato? \\
Requirement split & Le clausole formano una sola obbligazione coerente? Possono evolvere/essere assessed indipendentemente? \\
SR & Strengthening di un solo FR? Removal test? Niente duplicazione funzionale o realization leakage? \\
SecurityRequirement & Una sola security property? Failure mode identificabile? Niente generic ``secure'', cause o control prematuro? \\
Consumed service & Consumption o ownership? Se serve P, il service governa davvero P? NOT SPECIFIED resta NOT SPECIFIED? \\
Terminology & Stessa semantic identity, stesso governed term/binding? Marker temporanei rimossi dal final prose? \\
Propagation & La correzione upstream e' preservata nei child? Ci sono silent preconditions o semantic loss? \\
Promotion & Successor completo? Gap espliciti? Authority registry aggiornato? \\
BA handoff & Posso derivare la minimum BA senza inventare? Se la BA diverge, quale diagnosi e' plausibile? \\
\bottomrule
\end{longtable}

\section{Known open items -- non chiudere per comodita'}

\begin{statusopen}
La guida conserva esplicitamente i seguenti punti aperti o evidence-gated.
\end{statusopen}

\begin{enumerate}
\item rappresentazione document-layer/L1 esatta del target di un consumed service/capability;
\item universal service-property contract;
\item automatic residual-obligation generation;
\item provenance/change-event model completo fra analysis e governed-document revisions;
\item general role-aware BA topology regression method;
\item integrated Base Analysis / BA6 closure;
\item necessita'/ridondanza finale dei current BA2 operators dopo full-corpus retest;
\item eventuale reopen della regular Decision layer solo davanti a concrete recurring counterexamples;
\item universal taxonomy of semantic patterns;
\item network-specific metamodel.
\end{enumerate}

\section{Evidence map e status degli esempi}

Gli esempi non hanno tutti la stessa authority. La guida usa esplicitamente il loro status:

\begin{description}
\item[Forward checkpoint.] \nolinkurl{methodology/DDTA_DOCUMENTATION_METHOD_BASELINE_R24_CHECKPOINT_R1.md} -- question set e status epistemici correnti.
\item[FR authoring closure evidence.] \nolinkurl{03-functional-requirement/02-authoring-guidance/DDTA_AUTHORING_RULES_THROUGH_FR_R2.md} -- framing, MR, Decision, FR e service composition.
\item[S1 SpecializedRequirement closure.] CLOSED EVIDENCE -- strengthening, removal test e parent dependence.
\item[S1.5-A Requirement abstraction.] CLOSED nel proprio scope -- \code{Requirement}, \code{normativeClause[1..*]} e split rule.
\item[S2 SecurityRequirement closure.] CLOSED EVIDENCE -- \code{protectedSecurityProperty} e failure-mode semantics.
\item[MR-0003 R1.] \nolinkurl{governed-corpora/facial-access/current/MR-0003_IDENTITY_VERIFICATION.md} -- SUPERSEDED regression evidence; usato per mostrare l'ambiguita' pre-correction.
\item[MR-0003 candidate R2.] \nolinkurl{governed-corpora/facial-access/candidate-r2/MR-0003_IDENTITY_DETERMINATION.md} -- EXPERIMENTAL\_NON\_CANONICAL; usato come esempio di wording candidato corretto.
\item[Candidate Decisions.] \nolinkurl{governed-corpora/facial-access/candidate-r2/DECISIONS.md} -- EXPERIMENTAL\_NON\_CANONICAL; esempi D-3.1 e D-3.2.
\item[Candidate FunctionalRequirement.] \nolinkurl{governed-corpora/facial-access/candidate-r2/FUNCTIONAL_REQUIREMENTS.md} -- EXPERIMENTAL\_NON\_CANONICAL; esempio FR-3.2.1.
\item[SecurityRequirements R1.] \nolinkurl{governed-corpora/facial-access/current/SECURITY_REQUIREMENTS.md} -- SUPERSEDED regression evidence; esempi mechanism-neutral di confidentiality, integrity e authorized provenance.
\end{description}

\section{Final acceptance rule for this guide}

Una rappresentazione documentale e' accettabile quando e':

\begin{enumerate}
\item leggibile da autori e reviewer di progetto;
\item fedele al livello di astrazione governato;
\item semanticamente sufficiente per evitare letture materialmente divergenti dove il progetto ha gia' una risposta;
\item esplicitamente incompleta dove il progetto non ha ancora governato la risposta;
\item abbastanza testabile al livello FR/SR pertinente senza inventare meccanismi;
\item abbastanza stabile da poter evolvere la realization senza riscrivere responsabilita' che non sono cambiate;
\item pronta per una Base Analysis minima che normalizza, non crea, il significato.
\end{enumerate}

\begin{statusworking}[R24 DRAFT DISPOSITION]
Questa R2 e' una guida operativa derivata dal checkpoint documentale. Prima della promozione forward definitiva deve essere riletta contro il checkpoint, verificata sugli esempi Facial Access e poi usata per il successivo full documentation rewrite. Le regole R24 WORKING restano esplicitamente candidate finche' ulteriori pressure test non ne giustificano la closure.
\end{statusworking}

\end{document}
